
\documentclass{article}
\usepackage{colortbl}
\usepackage{makecell}
\usepackage{multirow}
\usepackage{supertabular}

\begin{document}

\newcounter{utterance}

\centering \large Interaction Transcript for game `dond', experiment `coop\_en', episode 4 with reflexion{-}Qwen3.5{-}9B.
\vspace{24pt}

{ \footnotesize  \setcounter{utterance}{1}
\setlength{\tabcolsep}{0pt}
\begin{supertabular}{c@{$\;$}|p{.15\linewidth}@{}p{.15\linewidth}p{.15\linewidth}p{.15\linewidth}p{.15\linewidth}p{.15\linewidth}}
    \# & $\;$A & \multicolumn{4}{c}{Game Master} & $\;\:$B\\
    \hline

    \theutterance \stepcounter{utterance}  
    & & \multicolumn{4}{p{0.6\linewidth}}{
        \cellcolor[rgb]{0.9,0.9,0.9}{
            \makecell[{{p{\linewidth}}}]{
                \texttt{\tiny{[P1$\langle$GM]}}
                \texttt{You have attempted this same task before and did not succeed. These are the plans you wrote after those attempts.} \\
\texttt{Trial \#1: New plan: I failed because I revealed my secret value function in the first message, which violated the rule that messages (unless secret proposals) are directly transmitted and should not include responses to rules or internal thoughts. Additionally, I used the square bracket format for a non{-}secret message, which the Game Master interpreted as a secret proposal, causing a conflict. Next time, I will strictly adhere to the format: I will send only the items I want in plain text without square brackets to communicate openly, and I will reserve the exact format "{[}Proposal: ...{]}" exclusively for my final secret submission. I will also ensure I do not accidentally paste my internal reasoning or value functions into the chat.} \\
\\ 
\\ 
\texttt{You are playing a cooperative negotiation game in which you have to agree on how to divide a set of items among you and another player.} \\
\\ 
\texttt{Rules:} \\
\texttt{(a) You and the other player are given a set of items. Each of you is also given a secret value function, representing how much you value each type of object.} \\
\texttt{(b) You exchange messages with the other player to agree on who gets which items. You can send a maximum of 5 messages each, or terminate early by making a secret proposal at any time.} \\
\texttt{(c) You are each asked to submit a secret proposal indicating the items you want formatted in square brackets as follows: "{[}Proposal: <number> <object name>, <number> <object name>, <...>{]}"} \\
\texttt{(d) If your proposals are complementary, i.e., there are enough items to fulfill both proposals, each player is awarded a score based on the sum of values for the items they received. Otherwise, both of you get zero points.} \\
\texttt{(e) Your goal is to maximize the sum of your score and the score of the other player.} \\
\\ 
\texttt{Let us start.} \\
\\ 
\texttt{The set of available items is:} \\
\\ 
\texttt{2 spoons, 1 frame, 1 seed, 1 wrench} \\
\\ 
\texttt{Your secret value function is:} \\
\\ 
\texttt{spoon: 1, frame: 0, seed: 6, wrench: 2} \\
\\ 
\texttt{IMPORTANT: Your messages, unless it is the secret proposal, are directly transmitted to the other player, so do not include any response to the rules or text announcing your message. To make a secret proposal, use the indicated format. Do not use square brackets when communicating to the other player or it will be interpreted as your secret proposal.} \\
            }
        }
    }
    & \\ \\

    \theutterance \stepcounter{utterance}  
    & \multicolumn{4}{p{0.6\linewidth}}{
        \cellcolor[rgb]{0.9,0.9,0.9}{
            \makecell[{{p{\linewidth}}}]{
                \texttt{\tiny{[P1$\rangle$GM]}}
                \texttt{Hello! I'm ready to negotiate. Let's work together to maximize our combined score.} \\
\\ 
\texttt{I see we have 2 spoons, 1 frame, 1 seed, and 1 wrench available. Based on my valuation, the seed is very valuable to me (6 points), while the wrench is worth 2 points, and the spoons are worth 1 point each. The frame isn't valuable to me at all.} \\
\\ 
\texttt{To ensure we both get high scores, I think we should focus on getting the seed and the wrench, as those are the high{-}value items for me. The spoons and frame can be distributed to you if they are valuable to you, or we can split them if needed.} \\
\\ 
\texttt{Let's start by discussing who might want the seed and the wrench. What are your valuations for these items?} \\
            }
        }
    }
    & & \\ \\

    \theutterance \stepcounter{utterance}  
    & & & \multicolumn{2}{p{0.3\linewidth}}{
        \cellcolor[rgb]{0.9,0.9,0.9}{
            \makecell[{{p{\linewidth}}}]{
                \texttt{\tiny{[GM$|$GM]}}
                \texttt{continue} \\
            }
        }
    }
    & & \\ \\

    \theutterance \stepcounter{utterance}  
    & & & \multicolumn{2}{p{0.3\linewidth}}{
        \cellcolor[rgb]{0.9,0.9,0.9}{
            \makecell[{{p{\linewidth}}}]{
                \texttt{\tiny{[GM$|$GM]}}
                \texttt{Hello! I'm ready to negotiate. Let's work together to maximize our combined score.} \\
\\ 
\texttt{I see we have 2 spoons, 1 frame, 1 seed, and 1 wrench available. Based on my valuation, the seed is very valuable to me (6 points), while the wrench is worth 2 points, and the spoons are worth 1 point each. The frame isn't valuable to me at all.} \\
\\ 
\texttt{To ensure we both get high scores, I think we should focus on getting the seed and the wrench, as those are the high{-}value items for me. The spoons and frame can be distributed to you if they are valuable to you, or we can split them if needed.} \\
\\ 
\texttt{Let's start by discussing who might want the seed and the wrench. What are your valuations for these items?} \\
            }
        }
    }
    & & \\ \\

    \theutterance \stepcounter{utterance}  
    & & \multicolumn{4}{p{0.6\linewidth}}{
        \cellcolor[rgb]{0.9,0.9,0.9}{
            \makecell[{{p{\linewidth}}}]{
                \texttt{\tiny{[GM$\rangle$P2]}}
                \texttt{You have attempted this same task before and did not succeed. These are the plans you wrote after those attempts.} \\
\texttt{Trial \#1: New plan: I failed because I included explanatory text and role clarification in my message, which violated the rule that non{-}proposal messages must be transmitted directly to the other player without any extra commentary. Additionally, I used square brackets in a non{-}proposal context, which the system interpreted as a secret proposal. Next time, I will strictly adhere to the format: if sending a public message, I will write only the raw text of the negotiation (e.g., "I propose we split the items as follows...") with absolutely no brackets, explanations, or meta{-}commentary. If I need to make a secret proposal, I will ensure it is the *only* content in the message and formatted exactly as `{[}Proposal: ...{]}`. I will start by sending a clear, bracket{-}free public message proposing a specific split based on our values.} \\
\\ 
\\ 
\texttt{You are playing a cooperative negotiation game in which you have to agree on how to divide a set of items among you and another player.} \\
\\ 
\texttt{Rules:} \\
\texttt{(a) You and the other player are given a set of items. Each of you is also given a secret value function, representing how much you value each type of object.} \\
\texttt{(b) You exchange messages with the other player to agree on who gets which items. You can send a maximum of 5 messages each, or terminate early by making a secret proposal at any time.} \\
\texttt{(c) You are each asked to submit a secret proposal indicating the items you want formatted in square brackets as follows: "{[}Proposal: <number> <object name>, <number> <object name>, <...>{]}"} \\
\texttt{(d) If your proposals are complementary, i.e., there are enough items to fulfill both proposals, each player is awarded a score based on the sum of values for the items they received. Otherwise, both of you get zero points.} \\
\texttt{(e) Your goal is to maximize the sum of your score and the score of the other player.} \\
\\ 
\texttt{Let us start.} \\
\\ 
\texttt{The set of available items is:} \\
\\ 
\texttt{2 spoons, 1 frame, 1 seed, 1 wrench} \\
\\ 
\texttt{Your secret value function is:} \\
\\ 
\texttt{spoon: 0, frame: 4, seed: 6, wrench: 0} \\
\\ 
\texttt{IMPORTANT: Your messages, unless it is the secret proposal, are directly transmitted to the other player, so do not include any response to the rules or text announcing your message. To make a secret proposal, use the indicated format. Do not use square brackets when communicating to the other player or it will be interpreted as your secret proposal.} \\
\\ 
\\ 
\texttt{Hello! I'm ready to negotiate. Let's work together to maximize our combined score.} \\
\\ 
\texttt{I see we have 2 spoons, 1 frame, 1 seed, and 1 wrench available. Based on my valuation, the seed is very valuable to me (6 points), while the wrench is worth 2 points, and the spoons are worth 1 point each. The frame isn't valuable to me at all.} \\
\\ 
\texttt{To ensure we both get high scores, I think we should focus on getting the seed and the wrench, as those are the high{-}value items for me. The spoons and frame can be distributed to you if they are valuable to you, or we can split them if needed.} \\
\\ 
\texttt{Let's start by discussing who might want the seed and the wrench. What are your valuations for these items?} \\
            }
        }
    }
    & \\ \\

    \theutterance \stepcounter{utterance}  
    & & & \multicolumn{4}{p{0.6\linewidth}}{
        \cellcolor[rgb]{0.9,0.9,0.9}{
            \makecell[{{p{\linewidth}}}]{
                \texttt{\tiny{[GM$\langle$P2]}}
                \texttt{I propose we give you the seed and I take the frame, while we split the spoons and wrench between us.} \\
            }
        }
    }
     \\ \\

    \theutterance \stepcounter{utterance}  
    & & & \multicolumn{2}{p{0.3\linewidth}}{
        \cellcolor[rgb]{0.9,0.9,0.9}{
            \makecell[{{p{\linewidth}}}]{
                \texttt{\tiny{[GM$|$GM]}}
                \texttt{continue} \\
            }
        }
    }
    & & \\ \\

    \theutterance \stepcounter{utterance}  
    & & & \multicolumn{2}{p{0.3\linewidth}}{
        \cellcolor[rgb]{0.9,0.9,0.9}{
            \makecell[{{p{\linewidth}}}]{
                \texttt{\tiny{[GM$|$GM]}}
                \texttt{I propose we give you the seed and I take the frame, while we split the spoons and wrench between us.} \\
            }
        }
    }
    & & \\ \\

    \theutterance \stepcounter{utterance}  
    & & \multicolumn{4}{p{0.6\linewidth}}{
        \cellcolor[rgb]{0.9,0.9,0.9}{
            \makecell[{{p{\linewidth}}}]{
                \texttt{\tiny{[P1$\langle$GM]}}
                \texttt{I propose we give you the seed and I take the frame, while we split the spoons and wrench between us.} \\
            }
        }
    }
    & \\ \\

    \theutterance \stepcounter{utterance}  
    & \multicolumn{4}{p{0.6\linewidth}}{
        \cellcolor[rgb]{0.9,0.9,0.9}{
            \makecell[{{p{\linewidth}}}]{
                \texttt{\tiny{[P1$\rangle$GM]}}
                \texttt{That sounds like a fair starting point. Since I value the seed at 6 points, getting it is a priority for me. The frame is worth 0 to me, so giving it to you is fine as long as you value it.} \\
\\ 
\texttt{Regarding the spoons and wrench: I value the wrench at 2 points and each spoon at 1 point. If you take the frame, I would prefer to take the wrench and one spoon, leaving you with the other spoon. This way, we both get items we value, and we avoid leaving anything on the table that could add to our combined score.} \\
\\ 
\texttt{Does that work for you? If so, I can make a secret proposal to finalize this split.} \\
            }
        }
    }
    & & \\ \\

    \theutterance \stepcounter{utterance}  
    & & & \multicolumn{2}{p{0.3\linewidth}}{
        \cellcolor[rgb]{0.9,0.9,0.9}{
            \makecell[{{p{\linewidth}}}]{
                \texttt{\tiny{[GM$|$GM]}}
                \texttt{continue} \\
            }
        }
    }
    & & \\ \\

    \theutterance \stepcounter{utterance}  
    & & & \multicolumn{2}{p{0.3\linewidth}}{
        \cellcolor[rgb]{0.9,0.9,0.9}{
            \makecell[{{p{\linewidth}}}]{
                \texttt{\tiny{[GM$|$GM]}}
                \texttt{That sounds like a fair starting point. Since I value the seed at 6 points, getting it is a priority for me. The frame is worth 0 to me, so giving it to you is fine as long as you value it.} \\
\\ 
\texttt{Regarding the spoons and wrench: I value the wrench at 2 points and each spoon at 1 point. If you take the frame, I would prefer to take the wrench and one spoon, leaving you with the other spoon. This way, we both get items we value, and we avoid leaving anything on the table that could add to our combined score.} \\
\\ 
\texttt{Does that work for you? If so, I can make a secret proposal to finalize this split.} \\
            }
        }
    }
    & & \\ \\

    \theutterance \stepcounter{utterance}  
    & & \multicolumn{4}{p{0.6\linewidth}}{
        \cellcolor[rgb]{0.9,0.9,0.9}{
            \makecell[{{p{\linewidth}}}]{
                \texttt{\tiny{[GM$\rangle$P2]}}
                \texttt{That sounds like a fair starting point. Since I value the seed at 6 points, getting it is a priority for me. The frame is worth 0 to me, so giving it to you is fine as long as you value it.} \\
\\ 
\texttt{Regarding the spoons and wrench: I value the wrench at 2 points and each spoon at 1 point. If you take the frame, I would prefer to take the wrench and one spoon, leaving you with the other spoon. This way, we both get items we value, and we avoid leaving anything on the table that could add to our combined score.} \\
\\ 
\texttt{Does that work for you? If so, I can make a secret proposal to finalize this split.} \\
            }
        }
    }
    & \\ \\

    \theutterance \stepcounter{utterance}  
    & & & \multicolumn{4}{p{0.6\linewidth}}{
        \cellcolor[rgb]{0.9,0.9,0.9}{
            \makecell[{{p{\linewidth}}}]{
                \texttt{\tiny{[GM$\langle$P2]}}
                \texttt{That works perfectly for me. I will take the frame (worth 4 to me) and one spoon (worth 1 to me), while you take the seed (worth 6 to you), the wrench (worth 2 to you), and the other spoon (worth 1 to you). This gives us a combined score of 14, which is the maximum possible since the total value of all items is 13 (6+4+2+1) plus the extra spoon value depending on distribution, but wait—let me re{-}calculate the total value.} \\
\\ 
\texttt{Total item values:} \\
\texttt{{-} Seed: 6 (me) / ? (you)} \\
\texttt{{-} Frame: 4 (me) / 0 (you)} \\
\texttt{{-} Wrench: 0 (me) / 2 (you)} \\
\texttt{{-} Spoons: 1 each (me) / 1 each (you)} \\
\\ 
\texttt{Wait, I need to check your valuations again. You said: "I value the seed at 6 points... wrench at 2 points and each spoon at 1 point." You didn't explicitly state your value for the frame, but implied it's low or zero since you're okay with me taking it. Let's assume your frame value is 0.} \\
\\ 
\texttt{My proposed split:} \\
\texttt{{-} Me: Frame (4) + 1 Spoon (1) = 5} \\
\texttt{{-} You: Seed (6) + Wrench (2) + 1 Spoon (1) = 9} \\
\texttt{{-} Total: 14} \\
\\ 
\texttt{But the total value of all items based on my values is: Seed(6) + Frame(4) + Wrench(0) + 2*Spoons(2) = 12.} \\
\texttt{Based on your values: Seed(6) + Frame(0) + Wrench(2) + 2*Spoons(2) = 10.} \\
\\ 
\texttt{The maximum combined score depends on who gets what. If I get the frame (4) and a spoon (1), and you get the seed (6), wrench (2), and a spoon (1), the sum is 4+1+6+2+1 = 14. But wait, the total value of items is fixed by the items themselves, not by who values them. The "combined score" is the sum of my valuation of my items plus your valuation of your items.} \\
\\ 
\texttt{Let's list all items and their values:} \\
\texttt{{-} Seed: Me=6, You=6 (you said "I value the seed at 6")} \\
\texttt{{-} Frame: Me=4, You=0 (implied)} \\
\texttt{{-} Wrench: Me=0, You=2} \\
\texttt{{-} Spoon 1: Me=1, You=1} \\
\texttt{{-} Spoon 2: Me=1, You=1} \\
\\ 
\texttt{Total possible combined score if we split perfectly:} \\
\texttt{If I get Frame (4) + Spoon (} \\
            }
        }
    }
     \\ \\

    \theutterance \stepcounter{utterance}  
    & & & \multicolumn{2}{p{0.3\linewidth}}{
        \cellcolor[rgb]{0.9,0.9,0.9}{
            \makecell[{{p{\linewidth}}}]{
                \texttt{\tiny{[GM$|$GM]}}
                \texttt{continue} \\
            }
        }
    }
    & & \\ \\

    \theutterance \stepcounter{utterance}  
    & & & \multicolumn{2}{p{0.3\linewidth}}{
        \cellcolor[rgb]{0.9,0.9,0.9}{
            \makecell[{{p{\linewidth}}}]{
                \texttt{\tiny{[GM$|$GM]}}
                \texttt{That works perfectly for me. I will take the frame (worth 4 to me) and one spoon (worth 1 to me), while you take the seed (worth 6 to you), the wrench (worth 2 to you), and the other spoon (worth 1 to you). This gives us a combined score of 14, which is the maximum possible since the total value of all items is 13 (6+4+2+1) plus the extra spoon value depending on distribution, but wait—let me re{-}calculate the total value.} \\
\\ 
\texttt{Total item values:} \\
\texttt{{-} Seed: 6 (me) / ? (you)} \\
\texttt{{-} Frame: 4 (me) / 0 (you)} \\
\texttt{{-} Wrench: 0 (me) / 2 (you)} \\
\texttt{{-} Spoons: 1 each (me) / 1 each (you)} \\
\\ 
\texttt{Wait, I need to check your valuations again. You said: "I value the seed at 6 points... wrench at 2 points and each spoon at 1 point." You didn't explicitly state your value for the frame, but implied it's low or zero since you're okay with me taking it. Let's assume your frame value is 0.} \\
\\ 
\texttt{My proposed split:} \\
\texttt{{-} Me: Frame (4) + 1 Spoon (1) = 5} \\
\texttt{{-} You: Seed (6) + Wrench (2) + 1 Spoon (1) = 9} \\
\texttt{{-} Total: 14} \\
\\ 
\texttt{But the total value of all items based on my values is: Seed(6) + Frame(4) + Wrench(0) + 2*Spoons(2) = 12.} \\
\texttt{Based on your values: Seed(6) + Frame(0) + Wrench(2) + 2*Spoons(2) = 10.} \\
\\ 
\texttt{The maximum combined score depends on who gets what. If I get the frame (4) and a spoon (1), and you get the seed (6), wrench (2), and a spoon (1), the sum is 4+1+6+2+1 = 14. But wait, the total value of items is fixed by the items themselves, not by who values them. The "combined score" is the sum of my valuation of my items plus your valuation of your items.} \\
\\ 
\texttt{Let's list all items and their values:} \\
\texttt{{-} Seed: Me=6, You=6 (you said "I value the seed at 6")} \\
\texttt{{-} Frame: Me=4, You=0 (implied)} \\
\texttt{{-} Wrench: Me=0, You=2} \\
\texttt{{-} Spoon 1: Me=1, You=1} \\
\texttt{{-} Spoon 2: Me=1, You=1} \\
\\ 
\texttt{Total possible combined score if we split perfectly:} \\
\texttt{If I get Frame (4) + Spoon (} \\
            }
        }
    }
    & & \\ \\

    \theutterance \stepcounter{utterance}  
    & & \multicolumn{4}{p{0.6\linewidth}}{
        \cellcolor[rgb]{0.9,0.9,0.9}{
            \makecell[{{p{\linewidth}}}]{
                \texttt{\tiny{[P1$\langle$GM]}}
                \texttt{That works perfectly for me. I will take the frame (worth 4 to me) and one spoon (worth 1 to me), while you take the seed (worth 6 to you), the wrench (worth 2 to you), and the other spoon (worth 1 to you). This gives us a combined score of 14, which is the maximum possible since the total value of all items is 13 (6+4+2+1) plus the extra spoon value depending on distribution, but wait—let me re{-}calculate the total value.} \\
\\ 
\texttt{Total item values:} \\
\texttt{{-} Seed: 6 (me) / ? (you)} \\
\texttt{{-} Frame: 4 (me) / 0 (you)} \\
\texttt{{-} Wrench: 0 (me) / 2 (you)} \\
\texttt{{-} Spoons: 1 each (me) / 1 each (you)} \\
\\ 
\texttt{Wait, I need to check your valuations again. You said: "I value the seed at 6 points... wrench at 2 points and each spoon at 1 point." You didn't explicitly state your value for the frame, but implied it's low or zero since you're okay with me taking it. Let's assume your frame value is 0.} \\
\\ 
\texttt{My proposed split:} \\
\texttt{{-} Me: Frame (4) + 1 Spoon (1) = 5} \\
\texttt{{-} You: Seed (6) + Wrench (2) + 1 Spoon (1) = 9} \\
\texttt{{-} Total: 14} \\
\\ 
\texttt{But the total value of all items based on my values is: Seed(6) + Frame(4) + Wrench(0) + 2*Spoons(2) = 12.} \\
\texttt{Based on your values: Seed(6) + Frame(0) + Wrench(2) + 2*Spoons(2) = 10.} \\
\\ 
\texttt{The maximum combined score depends on who gets what. If I get the frame (4) and a spoon (1), and you get the seed (6), wrench (2), and a spoon (1), the sum is 4+1+6+2+1 = 14. But wait, the total value of items is fixed by the items themselves, not by who values them. The "combined score" is the sum of my valuation of my items plus your valuation of your items.} \\
\\ 
\texttt{Let's list all items and their values:} \\
\texttt{{-} Seed: Me=6, You=6 (you said "I value the seed at 6")} \\
\texttt{{-} Frame: Me=4, You=0 (implied)} \\
\texttt{{-} Wrench: Me=0, You=2} \\
\texttt{{-} Spoon 1: Me=1, You=1} \\
\texttt{{-} Spoon 2: Me=1, You=1} \\
\\ 
\texttt{Total possible combined score if we split perfectly:} \\
\texttt{If I get Frame (4) + Spoon (} \\
            }
        }
    }
    & \\ \\

    \theutterance \stepcounter{utterance}  
    & \multicolumn{4}{p{0.6\linewidth}}{
        \cellcolor[rgb]{0.9,0.9,0.9}{
            \makecell[{{p{\linewidth}}}]{
                \texttt{\tiny{[P1$\rangle$GM]}}
                \texttt{That sounds like a perfect agreement! We have successfully identified a distribution that maximizes our combined score.} \\
\\ 
\texttt{To finalize this, I will now submit my secret proposal in the required format.} \\
\\ 
\texttt{{[}Proposal: 1 seed, 1 wrench, 1 spoon{]}} \\
            }
        }
    }
    & & \\ \\

    \theutterance \stepcounter{utterance}  
    & & & \multicolumn{2}{p{0.3\linewidth}}{
        \cellcolor[rgb]{0.9,0.9,0.9}{
            \makecell[{{p{\linewidth}}}]{
                \texttt{\tiny{[GM$|$GM]}}
                \texttt{proposal} \\
            }
        }
    }
    & & \\ \\

    \theutterance \stepcounter{utterance}  
    & & & \multicolumn{2}{p{0.3\linewidth}}{
        \cellcolor[rgb]{0.9,0.9,0.9}{
            \makecell[{{p{\linewidth}}}]{
                \texttt{\tiny{[GM$|$GM]}}
                [1, 0, 1, 1]
            }
        }
    }
    & & \\ \\

    \theutterance \stepcounter{utterance}  
    & & \multicolumn{4}{p{0.6\linewidth}}{
        \cellcolor[rgb]{0.9,0.9,0.9}{
            \makecell[{{p{\linewidth}}}]{
                \texttt{\tiny{[GM$\rangle$P2]}}
                \texttt{The other player has submitted a secret proposal. It is now your turn to submit a proposal of your own.} \\
            }
        }
    }
    & \\ \\

    \theutterance \stepcounter{utterance}  
    & & & \multicolumn{4}{p{0.6\linewidth}}{
        \cellcolor[rgb]{0.9,0.9,0.9}{
            \makecell[{{p{\linewidth}}}]{
                \texttt{\tiny{[GM$\langle$P2]}}
                \texttt{{[}Proposal: 1 frame, 1 spoon{]}} \\
            }
        }
    }
     \\ \\

    \theutterance \stepcounter{utterance}  
    & & & \multicolumn{2}{p{0.3\linewidth}}{
        \cellcolor[rgb]{0.9,0.9,0.9}{
            \makecell[{{p{\linewidth}}}]{
                \texttt{\tiny{[GM$|$GM]}}
                \texttt{proposal} \\
            }
        }
    }
    & & \\ \\

    \theutterance \stepcounter{utterance}  
    & & & \multicolumn{2}{p{0.3\linewidth}}{
        \cellcolor[rgb]{0.9,0.9,0.9}{
            \makecell[{{p{\linewidth}}}]{
                \texttt{\tiny{[GM$|$GM]}}
                [1, 1, 0, 0]
            }
        }
    }
    & & \\ \\

    \theutterance \stepcounter{utterance}  
    & & & \multicolumn{2}{p{0.3\linewidth}}{
        \cellcolor[rgb]{0.9,0.9,0.9}{
            \makecell[{{p{\linewidth}}}]{
                \texttt{\tiny{[GM$|$GM]}}
                [[1, 0, 1, 1], [1, 1, 0, 0]]
            }
        }
    }
    & & \\ \\

\end{supertabular}
}

\end{document}
